
\section{Black holes, wormholes, and recursion: channel management at finite bandwidth}
\label{sec:blackholes_recursion}

If physics is a readout, then black holes, singularities, and nonlocal phenomena can be restated as readout-theoretic boundary conditions.
This section formulates operational definitions compatible with the finite-resolution folding model, while keeping the minimal testable anchor separate from speculative extensions.

\subsection{Horizons as bandwidth saturation and computational boundaries}
\label{subsec:horizon_bandwidth}

In general relativity, horizons are geometric surfaces with strong causal significance \cite{Wald1984}.
In a readout cosmology, we interpret a horizon as a \emph{computational boundary} generated by saturation of local readout bandwidth and stability capacity.
This viewpoint is compatible with the area-law capacity paradigm and black-hole thermodynamics \cite{Bekenstein1973,Hawking1975}.

\begin{definition}[Horizon as channel saturation (operational)]
\label{def:horizon_saturation}
A region $\mathcal{R}$ is said to be \emph{horizon-saturated} if its effective input entropy density, combined with routing overhead, exceeds the capacity of the stable sector and the local routing budget.
The boundary of the maximal saturated region is the \emph{readout horizon}.
\end{definition}

The definition is intentionally operational: it makes horizon formation a consequence of finite capacity and stability constraints rather than a primitive geometric singularity.
In this view, ``firewall''-type behavior is reinterpreted as a protective mechanism for maintaining global readout consistency under local overload \cite{AlmheiriMarolfPolchinskiSully2013}.

\subsection{Singularities as resolution resets and recursive uplift}
\label{subsec:singularity_reset}

We treat ``singularity'' as an effective, non-geometric failure of a fixed-resolution readout interface.
Let $\chi(x)$ be the overhead proxy field induced by folding statistics (Definition~\ref{def:degeneracy_overhead}) and fix a threshold $\chi_{\mathrm{crit}}>0$.

\begin{definition}[Reset-critical region]
\label{def:reset_critical}
A region $\mathcal{R}$ is \emph{reset-critical} if $\sup_{x\in\mathcal{R}}\chi(x)\ge \chi_{\mathrm{crit}}$.
The boundary of the maximal reset-critical region is the \emph{reset boundary}.
\end{definition}

Operationally, a \emph{resolution reset} is a change in window length, address family, or internal phase register that re-encodes the local state into a new effective layer once the reset boundary is reached.

\begin{axiom}[Recursive reset interface]
\label{ax:recursive_reset}
There exists an iterability interface across resolution layers such that, beyond a saturation threshold, the readout switches to a new effective encoding while preserving the channel constraints as invariants of the rule set.
\end{axiom}

Under this axiom, a singularity is not an endpoint but a boundary between layers.
The internal discrete phase state (e.g.\ a $\ZZ_{2^p}$ register) can be treated as an initial hash for the subsequent layer: it seeds perturbations while inheriting the same channel rules.
This provides a natural model of recursive ``sub-universe'' branching without changing local laws.

\subsection{Wormholes as address pointer jumps}
\label{subsec:wormholes_pointer}

This subsection is an extension layer beyond the minimal $64\to 21$ model.
Hilbert addressing enforces locality by requiring that the scan traverses adjacent grid points.
A wormhole is defined as a controlled violation of this constraint: a nonlocal pointer jump in the address space.

\begin{definition}[Wormhole as an address pointer jump]
\label{def:wormhole_pointer_jump}
Let $H_n$ be a Hilbert address map at order $n$.
A \emph{wormhole link} is a directed (or undirected) pointer
\[
a \xrightarrow{\ \mathrm{ptr}\ } b,
\]
between scan indices $a,b\in\{0,\dots,2^{2n}-1\}$, implemented as a readout-level shortcut that bypasses the local adjacency traversal between $a$ and $b$.
\end{definition}

In routing language, wormholes trade physical locality for reduced dilation: they can reduce or neutralize scan cost at the expense of additional resources required to maintain channel consistency.
In the $\varphi$--$\pi$--$\e$ framework, a pointer jump risks violating grammar stability (unexpected local patterns), closure stability (monodromy), or analytic stability (pole barrier).
Maintaining a stable wormhole therefore requires additional overhead or a protected low-resolution tunnel, which is the readout-theoretic analog of the ``exotic matter'' requirement.



